\section{Introduction}
\label{sec:intro}

A dense 27-billion-parameter model quantized to about 4.5 bits per weight occupies roughly
17\,GB, which leaves about 6\,GB of a 24\,GB consumer GPU for the key--value (KV) cache, the
activations and the engine's own buffers. At batch size one, decoding is bound by memory
bandwidth: every generated token reads every weight once. On an RTX~4090 (1008\,GB/s), this
puts plain autoregressive decoding of such a model at about 50--60 tokens per second, whatever
the engine.

Two levers move this limit without changing the weights. \emph{Speculative decoding}
 lets a cheap drafter propose several tokens that the target model
verifies in one forward pass, which costs about as much memory traffic as a single token. And
\emph{engine specialization} removes the overheads a general engine carries: kernel launches,
host round trips and generic memory layouts. A third lever, \emph{KV-cache quantization}, does
not make short prompts faster, but decides how much context fits in the memory that is left.

This report studies the three levers on one concrete case: Qwen3.8-27B, a hybrid
Gated-DeltaNet/attention model with a built-in multi-token-prediction (MTP) head, served for a
single user who mostly writes code, on one RTX~4090. We start from a tuned llama.cpp setup and
ask three questions:

\begin{enumerate}
  \item How much faster can single-stream decoding get, at unchanged accuracy?
  \item How much context fits on the card, with and without the vision encoder, without the
        server failing?
  \item What does a 4-bit KV cache cost in precision on code, compared with the 8-bit keys
        and 4-bit values of the starting setup?
\end{enumerate}

We compare four engines: llama.cpp with the TurboQuant KV types \citep{llamacpp,turboquant}
(the baseline), the OpenJet persistent megakernel \citep{openjet}, the RTX~4090 port of NInfer
\citep{ninfer,ninfer4090} with both MTP and the DFlash2 block-diffusion drafter,
 and azllm \citep{azllm}, a fork of tqLLM \citep{tqllm} that reads
rotated TurboQuant weights. All measurements come from one machine, in throw-away containers,
with scripts that are part of this work.

\paragraph{Contributions.}
\begin{itemize}
  \item A like-for-like measurement of decode speed, accuracy and context capacity for four
        engines serving the same model on the same GPU, including thinking mode and a 29K-token
        prompt.
  \item A comparison of three 4-bit weight files on exactly the same code positions, showing
        that a rotated 4.3-bit format predicts code within 0.5\% of a 5.2-bit k-quant file.
  \item A capacity map: the largest context each configuration (engine, KV type, drafter,
        vision on or off) accepts, each verified by filling the context with real code and
        querying it.
  \item A measurement of the cost of KV-cache quantization on code perplexity at 64K and 128K
        context, and on exact recall of code at up to 252K tokens.
  \item Ready-to-run server scripts for every configuration; the recommended one serves
        2.5$\times$ faster than the baseline with 1.9$\times$ its context and vision, and the
        full 262K window with a 4-bit KV cache.
\end{itemize}
